% Source: https://artofproblemsolving.com/community/c3h1768109_altitude_of_an_orthodiagonal_right_trapezoid
% Stable identity: aops-topic-1768109-q01
Put $A=(0,0)$, $B=(a,0)$, $D=(0,h)$ and $C=(b,h)$, where $h=AD$.

\begin{asy}
usepackage("amsmath");
size(5.5cm);
defaultpen(fontsize(10pt));
real a=9,bb=4,h=6;
pair A=(0,0),B=(a,0),D=(0,h),C=(bb,h);
pair X=extension(A,C,B,D);
draw(A--B--C--D--cycle,linewidth(0.9));
draw(A--C,blue);draw(B--D,blue);
draw((0.4,0)--(0.4,0.4)--(0,0.4));draw((0,h-0.4)--(0.4,h-0.4)--(0.4,h));
pair u=unit(C-X),v=unit(B-X);draw(X+0.4*u--X+0.4*u+0.4*v--X+0.4*v,blue);
label("$A$",A,SW);label("$B$",B,SE);label("$C$",C,NE);label("$D$",D,NW);
label("$a$",(A+B)/2,S);label("$b$",(D+C)/2,N);label("$h$",(A+D)/2,W);
label("\mbox{}",(min(currentpicture,user=true)+max(currentpicture,user=true))/2,black); // keeps the text after the figure black when LaTeX typesets the labels inline
\end{asy}

Then $\overrightarrow{AC}=(b,h)$ and $\overrightarrow{BD}=(-a,h)$, and the diagonals are perpendicular exactly when
\[
\overrightarrow{AC}\cdot\overrightarrow{BD}=-ab+h^2=0 ,
\]
that is, $h=\sqrt{ab}$. $\blacksquare$

With similar triangles instead: $\angle ABD=90^\circ-\angle BAC$ (the diagonals meet at a right angle) and $\angle DAC=90^\circ-\angle BAC$ (since $\angle DAB=90^\circ$), so the right triangles $DAB$ and $CDA$ are similar, and $\dfrac{AD}{AB}=\dfrac{DC}{DA}$, i.e.\ $h^2=ab$. (The figure uses $a=9$, $b=4$, $h=6$.)
